\section{Boundary entropy from birth-target information}
\label{sec:area_law_birth_information}

\begin{theorem}[Boundary entropy from birth-target information]
    \label{thm:area_law_birth_information_boundary}
    \lean{TNLean.PEPS.AreaLaw.regionalEntropy_le_boundary_of_birth_information_estimates}
    \leanok
    \uses{def:area_law_ground_entropy,def:area_law_two_family_partition}
    Let $\Lambda\subset\mathbb Z^2$ be finite, let $q\ge1$ be an integer,
    and let $\Omega\in\bigotimes_{v\in\Lambda}\mathbb C^q$ be a unit vector.
    Suppose
    \begin{align*}
        A=D\mathbin{\dot\cup}\bigcup_iX_i\subseteq\Lambda
    \end{align*}
    is a finite ordered two-family partition, with labels
    $f_i\in\{0,1\}$. For each piece, put
    \begin{align*}
        P_i & =\bigcup_{j<i,\,f_j=f_i}X_j,
            & J_i                          & =(\Lambda\setminus A)\cup P_i.
    \end{align*}
    Let $Y_i\subseteq\Lambda$ be birth targets satisfying
    \begin{align*}
        X_i & \subseteq Y_i\subseteq A,
            & Y_i\cap P_i               & =\varnothing.
    \end{align*}
    The birth targets need not be pairwise disjoint.
    Let $n_i\in\mathbb N$, $c_D,c_S\in\mathbb R$ and $c_I\ge0$.
    Suppose that, for every $i$,
    \begin{align*}
        I_\Omega(Y_i:J_i) & \le c_I n_i^{-100}.
    \end{align*}
    Suppose also that
    \begin{align*}
        |D| & \le c_D|\partial_\Lambda A|,
            & \sum_i n_i^{-100}            & \le c_S|\partial_\Lambda A|.
    \end{align*}
    Here $I_\Omega(U:W)=S_\Omega(U)+S_\Omega(W)-S_\Omega(U\cup W)$
    is ordinary mutual information of the actual regional reductions.
    At a zero scale the weight is defined by $0^{-100}=0$;
    only positive scales occur in the manuscript application.
    Then
    \begin{align*}
        S_\Omega(A) & \le
        \left(c_D\log q+\frac{c_Ic_S}{2}\right)|\partial_\Lambda A|.
    \end{align*}
    This is the conditional implication in the proof of
    \cite[Theorem~1.1, lines~838--852 of Section~11]{OpenAI2026AreaLaw}.
    For that application $Y_i=A\cap T_i$. The construction of the partition
    and birth targets, the stated estimates, and the uniform choice of their
    constants remain separate requirements for the full area law.
\end{theorem}
\begin{proof}\leanok
    \uses{thm:area_law_boundary_assembly}
    Since $Y_i\subseteq A$ and $Y_i\cap P_i=\varnothing$, the birth
    target is disjoint from the entire subsystem $J_i$.
    Discarding the part $Y_i\setminus X_i$ of its first system therefore gives
    \begin{align*}
        I_\Omega(X_i:J_i) & \le I_\Omega(Y_i:J_i)
        \le c_I n_i^{-100}.
    \end{align*}
    Put $\varepsilon_i=c_I n_i^{-100}$. Because $c_I\ge0$,
    \begin{align*}
        \sum_i\varepsilon_i
         & =c_I\sum_i n_i^{-100}
        \le c_Ic_S|\partial_\Lambda A|.
    \end{align*}
    Apply Theorem~\ref{thm:area_law_boundary_assembly} with error coefficient
    $c_Ic_S$ and residual coefficient $c_D$. That estimate combines the
    two-family cancellation of Theorem~\ref{thm:area_law_two_family_entropy}
    with the regional dimension bound of
    Theorem~\ref{thm:area_law_entropy_dimension}, and gives the claimed
    coefficient.
\end{proof}

\begin{theorem}[Boundary entropy from neighborhood information]
    \label{thm:area_law_neighborhood_information_boundary}
    \lean{TNLean.PEPS.AreaLaw.regionalEntropy_le_boundary_of_neighborhood_information_estimates}
    \leanok
    \uses{def:area_law_ground_entropy,def:area_law_two_family_partition}
    Let $\Lambda\subset\mathbb Z^2$ be finite, let $q\ge1$ be an integer,
    and let $\Omega\in\bigotimes_{v\in\Lambda}\mathbb C^q$ be a unit vector.
    Let $A=D\mathbin{\dot\cup}\bigcup_iX_i$ be a finite ordered two-family
    partition, and set $P_i=\bigcup_{j<i,\,f_j=f_i}X_j$.
    Suppose that $X_i\subseteq Y_i\subseteq A$, $Y_i\subseteq N_i\subseteq\Lambda$,
    and $(\Lambda\setminus A)\cup P_i\subseteq\Lambda\setminus N_i$ for every $i$.
    Let $n_i\in\mathbb N$, $c_D,c_S\in\mathbb R$ and $c_I\ge0$.
    Suppose that the information estimate holds for every $i$, and that the
    residual and scale sum satisfy the following bounds:
    \begin{align*}
        I_\Omega(Y_i:\Lambda\setminus N_i) & \le c_I n_i^{-100},                                         \\
        |D|                                & \le c_D|\partial_\Lambda A|,
                                           & \sum_i n_i^{-100}            & \le c_S|\partial_\Lambda A|.
    \end{align*}
    Then
    \begin{align*}
        S_\Omega(A) & \le
        \left(c_D\log q+\frac{c_Ic_S}{2}\right)|\partial_\Lambda A|.
    \end{align*}
    The zero-scale convention is $0^{-100}=0$; the manuscript application
    uses positive scales. This is an auxiliary implication in
    \cite[Theorem~1.1, lines~815--852 of Section~11]{OpenAI2026AreaLaw}.
    The construction of the partition, the whole-union exclusion and the
    information estimates, with uniformly chosen constants, remain separate
    requirements for the full area law.
\end{theorem}
\begin{proof}\leanok
    \uses{thm:area_law_birth_information_boundary}
    Put $J_i=(\Lambda\setminus A)\cup P_i$.
    Since $Y_i\subseteq N_i$ and $J_i\subseteq\Lambda\setminus N_i$,
    $Y_i$ is disjoint from $J_i$ and in particular from $P_i$.
    Discarding part of the second system gives
    \begin{align*}
        I_\Omega(Y_i:J_i) & \le I_\Omega(Y_i:\Lambda\setminus N_i)
        \le c_I n_i^{-100}.
    \end{align*}
    Apply Theorem~\ref{thm:area_law_birth_information_boundary} to the same
    partition, birth targets and original vector $\Omega$.
\end{proof}
